% =====================================================================
% Method section (Lag-RLS / DARO). Requires amsmath, algorithm,
% algpseudocode and the acronym package (\ac{SOC}, \ac{SOH}, \ac{QoS},
% \ac{CPU}, \ac{DVFS}, \ac{RLS}, \ac{ECM}, \ac{OCV}, \ac{DOD}).
% Numbers in comments refer to the implementation in methods/*.py.
% Structure follows Fig. lag_rls_structure:
%   A. Problem Formulation  B. Real-Time Loop (estimator, decision-making)
%   C. Lagrangian Planner (safe energy budget, Lagrangian optimization)
% =====================================================================

\section{Methodology}
\label{sec:method}

Fig.~\ref{fig:lag_rls_structure} shows the proposed controller. It jointly selects the
driving speed and the \ac{CPU} frequency with two timescales.
The \emph{real-time loop} runs every control step: an estimator adapts the power model and
estimates the battery loss factor, the complexity level and the battery state (\ac{SOC} and
\ac{SOH}) from voltage, current and throughput measurements, and the decision-making selects the speed
and the \ac{DVFS} level that minimize a cost per meter.
The \emph{Lagrangian planner} runs every 10~steps: it computes the safe energy budget from the
estimated \ac{SOC} and \ac{SOH} and selects the energy price $\lambda^\ast$ used by the decision-making.
We first formulate the problem (Section~\ref{sec:method:problem}) and then describe the
real-time loop (Section~\ref{sec:method:rt}) and the planner (Section~\ref{sec:method:planner}).

\begin{figure*}[t]
  \centering
  \includegraphics[width=\textwidth]{figures/lag_rls_structure.pdf}
  \caption{Structure of the proposed controller. In the real-time loop (every 1~s), the
  estimator updates the \ac{SOH} and \ac{SOC}, the battery loss factor $\hat\eta$, the power model
  $\hat P$, and the complexity level $\hat c$ from the rover observations, and the decision-making selects the speed command $v^\ast$ and
  the \ac{DVFS} level $f^\ast$ that minimize the cost per meter, including the expected \ac{QoS}
  penalty under the complexity variation; the rover tracks $v^\ast$ subject to its acceleration
  limit, and $v$ and $f$ denote the measured speed and frequency. The Lagrangian planner (every 10~s) computes the safe
  energy budget from $\widehat{\mathrm{SOC}}$ and $\widehat{\mathrm{SOH}}$, predicts the remaining
  energy with $\hat\eta$, $\hat P$, and $\hat c$, and selects the energy
  price $\lambda^\ast$ that trades the mission time against the battery energy and, for
  DARO, its degradation; without the degradation terms ($w_D = 0$) the controller reduces
  to Lag-RLS.}
  \label{fig:lag_rls_structure}
\end{figure*}

% =====================================================================
\subsection{Problem Formulation}
\label{sec:method:problem}

At each control step $t$ (period $\Delta t = 1$~s) the controller chooses the action
$u_t = (v_t, f_t)$, where $v_t \in [v_{\min}, v_{\max}]$ is the speed command and
$f_t \in \mathcal{F}$ is one of the $|\mathcal{F}| = 11$ \ac{DVFS} levels (0.50--2.04~GHz).
The rover must cover a mission of length $L$ while the perception pipeline sustains the
throughput $q \geq q_{\mathrm{ref}}$; $q$ depends on the speed, the \ac{CPU} frequency and the environmental complexity $c$.
The complexity is determined by the environment the rover traverses, e.g., the number of events
the perception pipeline must process. It follows the complexity level $\bar c$ of the current
terrain and fluctuates around it from step to step as the scene changes,
$c_t = \bar c\,u_t$ with $u_t \sim \mathcal{U}(1-\delta, 1+\delta)$; neither $\bar c$ nor $c_t$ is
known to the controller. Since $c_t$ is random, the requirement is stated as a chance constraint: the throughput may fall below
$q_{\mathrm{ref}}$ in at most a fraction $\rho = 5\%$ of the steps.

\paragraph{Power and throughput models}
The electrical load power $P$ is the sum of the motor power, which depends on the
speed, and the computing and auxiliary power $P_{\mathrm{c}}$, which depends on the speed, the \ac{CPU}
frequency and the complexity~\cite{REF}.
The throughput model
$q = A - B s^{p} + (G + D s^{r})(f - f_{\mathrm{ref}})$, with the processing load
$s = v c / c_{\mathrm{ref}}$, decreases with the speed and the complexity and increases linearly with $f$.
A lower \ac{DVFS} level therefore saves computing power but raises the probability of a \ac{QoS} shortfall
when the complexity fluctuates. Because the model is linear in $f$, it can also be inverted for the complexity,
\begin{equation}
  c = q^{-1}(v, f, q),
  \label{eq:fq}
\end{equation}
which the estimator uses in Section~\ref{sec:method:estimation}.

\paragraph{Battery model}
The load is supplied by a lithium-polymer pack, modeled by a second-order \ac{ECM}
with an ohmic resistance $R_0$ and two RC branches $(R_1, C_1)$ and $(R_2, C_2)$. For a
constant load power $P_t$ over the step, the terminal voltage and the current are
\begin{equation}
  V_t = V_{\mathrm{oc}}(\mathrm{SOC}_t, \mathrm{SOH}) - v_{1,t} - v_{2,t} - R_0 I_t,
  \qquad P_t = V_t I_t,
  \label{eq:ecm}
\end{equation}
so that $I_t = 2P_t / \big(u_t + \sqrt{u_t^2 - 4 R_0 P_t}\big)$ with
$u_t = V_{\mathrm{oc}} - v_{1,t} - v_{2,t}$; a power demand with $u_t^2 < 4R_0P_t$ cannot be
delivered by the battery. The polarization voltages and the \ac{SOC} evolve as
\begin{equation}
  v_{j,t+1} = a_j\, v_{j,t} + (1 - a_j)\, R_j I_t,\;\; a_j = e^{-\Delta t/\tau_j},
  \qquad
  \mathrm{SOC}_{t+1} = \mathrm{SOC}_t - \frac{I_t\,\Delta t}{Q_n\,\mathrm{SOH}},
  \label{eq:soc}
\end{equation}
for $j \in \{1, 2\}$, where $Q_n$ is the rated capacity and $\tau_j = R_j C_j$ are fixed time
constants. Aging enters the model in three places: the usable capacity is $Q_n\,\mathrm{SOH}$, the
resistances increase as
\begin{equation}
  R_j(\mathrm{SOC}, \mathrm{SOH}) = \exp\!\big(p_j(\mathrm{SOC})\big)\,\big[1 + \kappa_j (1 - \mathrm{SOH})\big],
  \quad j \in \{0, 1, 2\},
  \label{eq:ecm_param}
\end{equation}
with $p_j$ a second-order polynomial and $\kappa_j$ an aging coefficient, and the \ac{OCV} curve
changes shape. The latter is interpolated between a fresh and an aged curve,
$V_{\mathrm{oc}}(\mathrm{SOC}, \mathrm{SOH}) = (1 - w)\,V^{\mathrm{fr}}_{\mathrm{oc}}(\mathrm{SOC}) + w\,V^{\mathrm{ag}}_{\mathrm{oc}}(\mathrm{SOC})$
with $w = \min\!\big(1, (1 - \mathrm{SOH}) / (1 - \mathrm{SOH}_{\mathrm{ag}})\big)$, each fitted
to measured pack discharge curves at $\mathrm{SOH} = 1$ and $\mathrm{SOH}_{\mathrm{ag}} = 0.85$.
Because the aged curve is lower at medium and low \ac{SOC}, an aged battery delivers less energy
per unit of charge, in addition to its smaller capacity. The pack temperature $T$ follows a lumped
thermal model,
$C_{\mathrm{th}}\,\dot T = R_0 I^2 + v_1^2/R_1 + v_2^2/R_2 - G_{\mathrm{th}}\,(T - T_{\mathrm{amb}})$.
The parameters are identified in Section~\ref{sec:exp:setup} (Table~\ref{tb:battery_parameter}).

\paragraph{Degradation model}
We describe the capacity fade by the semi-empirical cycle- and calendar-aging model of
Xu et al.~\cite{xu2016modeling}. Accordingly, the battery \ac{SOH} is expressed as
\begin{equation}
  \mathrm{SOH}(t) = \frac{Q^{\mathrm{ag}}_{\mathrm{act}}(t)}{Q_n}
  = \mathrm{SOH}_0 \exp\!\big(-f_d(t)\big),
  \label{eq:battery_soh_define}
\end{equation}
where $Q^{\mathrm{ag}}_{\mathrm{act}}$ is the actual capacity of the aged battery, $\mathrm{SOH}_0$
denotes the initial battery health, and $f_d(t)$ represents the accumulated degradation, the sum of
the cycle and the calendar aging:
\begin{equation}
  f_d(t) = k_{\mathrm{cyc}} \sum_{i=1}^{N(t)} n_i\, \widetilde{S}_{\delta_i} S_{\sigma_i} S_{T_i}
         + f_{\mathrm{cal}}(t),
  \label{eq:battery_fd}
\end{equation}
where $N(t)$ is the number of cycles identified up to time $t$, $k_{\mathrm{cyc}}$ is the reference
cycle-aging coefficient, and $n_i$ is the cycle-counting weight, equal to 1 for a full cycle and
0.5 for a half cycle. The terms $\delta_i$, $\sigma_i$, and $T_i$ represent the \ac{DOD}, average
\ac{SOC}, and average temperature of cycle $i$, respectively. The corresponding stress factors are
defined as
\begin{subequations}\label{eq:battery_soh_factor}
\begin{align}
  \widetilde{S}_{\delta_i} &= \frac{k_{\delta_1} + k_{\delta_3}}{k_{\delta_1}\delta_i^{k_{\delta_2}} + k_{\delta_3}}, \\
  S_{\sigma_i} &= e^{k_\sigma(\sigma_i - \sigma_{\mathrm{ref}})}, \\
  S_{T_i} &= e^{k_T (T_i - T_{\mathrm{ref}}) \frac{T_{\mathrm{ref}}}{T_i}},
\end{align}
\end{subequations}
where $k_\sigma$, $k_T$, $k_{\delta_1}$, $k_{\delta_2}$, and $k_{\delta_3}$ are empirical stress-factor
coefficients, while $\sigma_{\mathrm{ref}}$ and $T_{\mathrm{ref}}$ denote the reference \ac{SOC} and
absolute temperature, respectively. The \ac{DOD} stress factor is normalized to unity at a full
cycle, so that $k_{\mathrm{cyc}}$ is the capacity-fade rate of one full cycle at
$\sigma_{\mathrm{ref}}$ and $T_{\mathrm{ref}}$.
The calendar aging accrues with time,
\begin{equation}
  f_{\mathrm{cal}}(t) = k_t \int_0^{t} S_\sigma\big(\mathrm{SOC}(\tau)\big)\, S_T\big(T(\tau)\big)\, d\tau,
  \label{eq:xu_cal}
\end{equation}
with $S_\sigma$ and $S_T$ as in \eqref{eq:battery_soh_factor} and the calendar-aging coefficient $k_t$.
For a linear \ac{SOC} change from $\sigma_a$ to $\sigma_b$ over a period $t_p$ at a constant
temperature $T$, \eqref{eq:xu_cal} has the closed form
\begin{equation}
  \Delta f_{\mathrm{cal}}(\sigma_a, \sigma_b, t_p, T)
  = k_t\, t_p\, S_T\, \frac{S_\sigma(\sigma_b) - S_\sigma(\sigma_a)}{k_\sigma(\sigma_b - \sigma_a)},
  \label{eq:xu_cal_lin}
\end{equation}
which tends to $k_t\, t_p\, S_T\, S_\sigma(\sigma_a)$ for $\sigma_b \to \sigma_a$.

\paragraph{Degradation per mission}
Each mission consists of two half cycles ($n_i = 0.5$): the discharge from $\mathrm{SOC}_0$ to
$\mathrm{SOC}_{\mathrm{end}}$ during the mission time $T_m$, and the subsequent constant-current
recharge to $\mathrm{SOC}_0$ at the rate $C_{\mathrm{ch}} = 1$~C and the ambient temperature
$T_{\mathrm{amb}}$. Both half cycles have the \ac{DOD} and average \ac{SOC}
\begin{equation}
  \delta = \mathrm{SOC}_0 - \mathrm{SOC}_{\mathrm{end}},\qquad
  \bar\sigma = \tfrac{1}{2}(\mathrm{SOC}_0 + \mathrm{SOC}_{\mathrm{end}}),
  \label{eq:cycle_stats}
\end{equation}
and the discharge half cycle has the average temperature
$\bar T = \frac{1}{T_m}\int_0^{T_m} T(\tau)\, d\tau$. By \eqref{eq:battery_fd}, the degradation of
the mission is
\begin{equation}
\begin{aligned}
  \Delta f_d = {} & \tfrac{1}{2} k_{\mathrm{cyc}}\, \widetilde{S}_{\delta} S_{\bar\sigma}
      \big(S_{\bar T} + S_{T_{\mathrm{amb}}}\big) \\
    & + \sum_{t=0}^{T_m/\Delta t - 1} \Delta f_{\mathrm{cal}}(\mathrm{SOC}_t, \mathrm{SOC}_{t+1}, \Delta t, T_t) \\
    & + \Delta f_{\mathrm{cal}}\big(\mathrm{SOC}_{\mathrm{end}}, \mathrm{SOC}_0, t_{\mathrm{ch}}, T_{\mathrm{amb}}\big),
\end{aligned}
  \label{eq:fd_mission}
\end{equation}
where the first line is the cycle aging of the two half cycles, the second line is the calendar
aging along the mission trajectory, and the third line is the calendar aging during the charging
time $t_{\mathrm{ch}} = \delta / C_{\mathrm{ch}}$. By \eqref{eq:battery_soh_define}, the \ac{SOH}
after the mission is $\mathrm{SOH}\,e^{-\Delta f_d}$, and the capacity loss of the mission is
\begin{equation}
  \Delta\mathrm{SOH} = \mathrm{SOH}\,\big(1 - e^{-\Delta f_d}\big).
  \label{eq:dsoh}
\end{equation}
Since $\Delta\mathrm{SOH}$ is of the order of $10^{-5}$ per mission, the \ac{SOH} is treated as
constant within a mission. The controller influences $\Delta\mathrm{SOH}$ through the \ac{DOD},
which grows with the cell energy drawn, through the average \ac{SOC}, through the mission time in the
calendar term, and, via the Joule losses of the \ac{ECM}, through the temperature.

\paragraph{Lifetime}
Repeating a mission from $\mathrm{SOH}_0 = 1$ until the end of life $\mathrm{SOH}_{\mathrm{EOL}}$
accumulates the degradation $f_d$ from 0 to $f_{d,\mathrm{EOL}} = -\ln \mathrm{SOH}_{\mathrm{EOL}}$.
Because the degradation per mission $\Delta f_d$ changes as the battery ages (the controller adapts
to the \ac{SOH}, and the aged \ac{ECM} draws more charge for the same load energy), the number of
missions to the end of life is
\begin{equation}
  N_{\mathrm{EOL}} = \int_0^{f_{d,\mathrm{EOL}}} \frac{df_d}{\Delta f_d\big(\mathrm{SOH} = e^{-f_d}\big)},
  \label{eq:n_eol}
\end{equation}
and the distance to the end of life is $N_{\mathrm{EOL}} L$. We evaluate \eqref{eq:n_eol} by the
trapezoidal rule, with $\Delta f_d$ interpolated between missions simulated at several \ac{SOH} levels.

\paragraph{Objective}
The mission objective trades travel time, energy, \ac{QoS} shortfalls and the terminal \ac{SOC} reserve:
\begin{equation}
\begin{aligned}
  \min_{\{u_t\}} \;
  J = {} & w_T \frac{T_m}{T_{\mathrm{ref}}}
    + w_E \frac{E}{L}
    + \frac{w_Q}{T_{\mathrm{ref}}} \int_0^{T_m} \psi_t\, dt
    + w_R \,\phi\big(\mathrm{SOC}_{\mathrm{end}}\big) \\
  \text{s.t.} \quad &
  \Pr\big[q_t < q_{\mathrm{ref}}\big] \leq \rho,\;\;
  \mathrm{SOC}_t \geq \mathrm{SOC}_{\mathrm{cut}},
\end{aligned}
  \label{eq:objective}
\end{equation}
where $T_m$ is the mission time, $T_{\mathrm{ref}} = L / v_{\mathrm{ref}}$, $E = \int P\,dt$ is the
load energy, $\psi_t = \mathbb{1}[q_t < q_{\mathrm{ref}}] + [1 - q_t/q_{\mathrm{ref}}]^+$ penalizes
each \ac{QoS} shortfall by its occurrence and its depth, $\mathrm{SOC}_{\mathrm{cut}}$ is the cutoff \ac{SOC} at which the battery
is considered depleted and the mission fails, and
$\phi(\mathrm{SOC}_{\mathrm{end}}) = \min\!\big(1, \max(0, \mathrm{SOC}_{\mathrm{res}} - \mathrm{SOC}_{\mathrm{end}}) / \mathrm{SOC}_{\mathrm{res}}\big)$
penalizes ending the mission below the reserve $\mathrm{SOC}_{\mathrm{res}} = 5\%$.
Writing time and energy as integrals over the travelled distance $s$,
$T_m = \int_0^L ds / v$ and $E = \int_0^L e\,ds$ with the energy per meter
$e = P / v$, makes the problem a distance-domain optimal control problem.

\paragraph{Degradation-aware objective}
The objective \eqref{eq:objective} only covers the current mission. To account for the
battery's long-term cost, DARO adds the predicted capacity loss of the mission,
\begin{equation}
  J_{\mathrm{DA}} = J + w_D \,\frac{\Delta\mathrm{SOH}}{1 - \mathrm{SOH}_{\mathrm{EOL}}},
  \label{eq:da}
\end{equation}
where $w_D$ is the cost of one full battery life (from \ac{SOH}~=~1 to the end-of-life
$\mathrm{SOH}_{\mathrm{EOL}}$) and $\Delta\mathrm{SOH}$ is the capacity loss \eqref{eq:dsoh} of the
mission and the subsequent recharge. For $w_D = 0$, DARO reduces to Lag-RLS.

\paragraph{Decomposition}
The \ac{SOC} constraint couples all steps of the mission, whereas the \ac{QoS} constraint and the
remaining terms are local in distance. We relax the \ac{SOC} constraint with a multiplier
$\lambda \geq 0$, which acts as the price of battery energy. For a constant complexity, the
Hamiltonian of \eqref{eq:objective} then decouples along the distance, and the optimal action at every
point minimizes a cost per meter that depends on the mission only through $\lambda$. This
splits the controller into a real-time loop that minimizes the cost per meter at a given price
(Section~\ref{sec:method:rt}) and a planner that selects the price from the battery state
(Section~\ref{sec:method:planner}). This is the structure of the equivalent consumption
minimization strategy.

% =====================================================================
\subsection{Real-Time Loop}
\label{sec:method:rt}

At every control step, the estimator updates the models and the battery state from the
observations, and the decision-making selects $(v^\ast, f^\ast)$ at the current price $\lambda^\ast$.

% ---------------------------------------------------------------------
\subsubsection{Estimator}
\label{sec:method:estimation}

\paragraph{SOH and SOC}
The coulomb-counted \ac{SOC} is referenced to the rated capacity, so on an aged battery it
underestimates the charge drawn per unit \ac{SOC}. For each candidate $h$ on a grid
$\mathcal{H} = \{0.70, 0.71, \dots, 1.00\}$ we compute the corrected
$\mathrm{SOC}_t(h) = \mathrm{SOC}_0 - (\mathrm{SOC}_0 - \mathrm{SOC}_t^{\mathrm{cc}})/h$ and the
predicted steady-state terminal voltage of the \ac{ECM} \eqref{eq:ecm},
$\hat V_t(h) = V_{\mathrm{oc}}(\mathrm{SOC}_t(h), h) - I_t\big(R_0 + R_1 + R_2\big)$,
with the resistances $R_j(\mathrm{SOC}_t(h), h)$ of \eqref{eq:ecm_param} and the nominal aging coefficients
$\kappa_{0,1,2} = 1.25, 0.89, 4.27$ (Table~\ref{tb:battery_parameter}). The estimate
$\widehat{\mathrm{SOH}} = \arg\min_h \overline{(\hat V_t(h) - V_t)^2}$ minimizes an exponential
moving average of the squared residuals over near-steady samples
($|\Delta I| < 0.05$~A), and $\widehat{\mathrm{SOC}} = \mathrm{SOC}_t(\widehat{\mathrm{SOH}})$.
Both enter the safe energy budget \eqref{eq:eav}.

\paragraph{Loss factor}
The load energy $\int V I\, dt$ differs from the cell energy drawn by the internal losses.
The ratio
$\hat\eta = Q_n \widehat{\mathrm{SOH}} \int_{\mathrm{SOC}_t}^{\mathrm{SOC}_0} V_{\mathrm{oc}}\, d\sigma \,/\, \int_0^t V I\, d\tau$,
clipped to $[0.9, 1.3]$, converts predicted load energy into cell energy in \eqref{eq:lambda}.

\paragraph{Power model}
The nominal power models are identified on a reference platform and can deviate on the
deployed rover (payload, terrain, compute load). We correct them on-line by \ac{RLS} on the
measured load power $y_t = V_t I_t$:
\begin{equation}
  y_t = \boldsymbol\theta^\top \boldsymbol\varphi_t,\qquad
  \boldsymbol\varphi_t = \big[\, k_{\mathrm{th}}(v)(b_0 + b_1 v),\;
                          k_{\mathrm{th}}(v)\, b_3 v^3,\;
                          P_{\mathrm{c}}^{\mathrm{nom}} \,\big]^\top,
  \label{eq:rls}
\end{equation}
where $b_i$ are the nominal motor coefficients and $k_{\mathrm{th}}$ is the steady-state
motor thermal-loss factor. $\boldsymbol\theta$ is updated with a forgetting factor
$\rho = 0.9995$, projected onto $[0.3, 3]$, and only at steady speed ($|\Delta v| < 0.02$~m/s),
since acceleration is not modeled. The corrected model $\hat P = \boldsymbol\theta^\top \boldsymbol\varphi$
is used by both the decision-making and the planner. To make the $v^3$ term identifiable, the
scheduled speed is scaled by $\{1, 0.75, 1.15\}$ for 30~s each at the mission start.

\paragraph{Complexity}
The controller does not observe the complexity directly; it is recovered from the measured
throughput by \eqref{eq:fq},
$c_t = q^{-1}(v_t, f_t, q_t)$. To avoid reacting to the step-to-step variation, the estimate $\hat c$ of the
complexity level $\bar c$ is the mean of $c_t$ over the current segment. A new segment starts at once when $c_t$ exceeds the level by
25\% (to protect \ac{QoS}) and after three consecutive samples when it falls 25\% below it.

% ---------------------------------------------------------------------
\subsubsection{Decision-Making}
\label{sec:method:schedule}

Given the price $\lambda^\ast$ from the planner, the decision-making minimizes the cost per meter
\begin{equation}
  \ell_\lambda(v, f) = \frac{1}{v}\left[
    \Big(\frac{(1+\lambda)\,w_E}{L} + \hat\eta\,\mu_E\Big) \hat P
    + \frac{w_T}{T_{\mathrm{ref}}} + \mu_t
    + \frac{w_Q}{T_{\mathrm{ref}}}\, \bar\psi(v, f, \hat c) \right],
  \label{eq:lagrangian}
\end{equation}
where $\mu_E$ (per joule of cell energy) and $\mu_t$ (per second) are the degradation prices
provided by the planner (Section~\ref{sec:method:planner}); for Lag-RLS, $\mu_E = \mu_t = 0$.
$\lambda = 0$ gives the unconstrained time--energy trade-off, and a larger $\lambda$ shifts the
operating point towards the energy-optimal speed.

Instead of a fixed \ac{DVFS} safety margin, the decision-making prices the \ac{QoS} risk
with the expected penalty
\begin{equation}
  \bar\psi(v, f, \hat c) = \mathbb{E}_{c}\Big[\mathbb{1}[q < q_{\mathrm{ref}}]
    + \big[1 - q/q_{\mathrm{ref}}\big]^+\Big],
  \quad c = \hat c\, u\, (1+\varepsilon),
  \label{eq:qos_cost}
\end{equation}
where $u \sim \mathcal{U}(1-\delta, 1+\delta)$ is the known step-to-step variation and
$\varepsilon \sim \mathcal{U}(-0.03, 0.03)$ accounts for the error of the complexity estimate $\hat c$;
the expectation is evaluated by quadrature (41~$\times$~7 points).
The decision-making searches the speed grid $\mathcal{V}$ (161 points) and the \ac{DVFS} levels jointly:
\begin{equation}
  (v^\ast, f^\ast) = \arg\min_{v \in \mathcal{V},\, f \in \mathcal{F}}
           \ell_{\lambda^\ast}(v, f)
  \quad \text{s.t.} \quad \Pr\big[q < q_{\mathrm{ref}}\big] \leq \rho .
  \label{eq:schedule}
\end{equation}
A level that occasionally misses the requirement is therefore chosen only where the energy it
saves outweighs the expected penalty; without variation ($\delta = \varepsilon = 0$),
\eqref{eq:schedule} reduces to the lowest \ac{DVFS} level that meets the requirement.
The rover tracks the speed command $v^\ast$ subject to its acceleration limit, while $f^\ast$ is
applied immediately.

% =====================================================================
\subsection{Lagrangian Planner}
\label{sec:method:planner}
\label{sec:method:pricing}

Every $T_p = 10$~s the planner updates the price $\lambda^\ast$ from the estimated battery state.

\paragraph{Safe energy budget}
The energy the battery can deliver above a floor $\mathrm{SOC}_{\mathrm{f}}$ depends on both battery states:
\begin{equation}
  E_{\mathrm{av}}(\mathrm{SOC}, \mathrm{SOH}; \mathrm{SOC}_{\mathrm{f}})
  = Q_n \,\mathrm{SOH} \int_{\mathrm{SOC}_{\mathrm{f}}}^{\mathrm{SOC}}
    V_{\mathrm{oc}}(\sigma, \mathrm{SOH})\, d\sigma,
  \label{eq:eav}
\end{equation}
where $Q_n$ is the rated capacity and $V_{\mathrm{oc}}$ the \ac{OCV} curve, which also
changes with aging. The \ac{SOC} thus sets the short-term budget, and the \ac{SOH} scales
both the capacity and the voltage at which the charge is delivered.
The planner evaluates \eqref{eq:eav} at $(\widehat{\mathrm{SOC}}, \widehat{\mathrm{SOH}})$ for two floors:
the safe floor $\mathrm{SOC}_{\mathrm{f}}$, which is the cutoff \ac{SOC} plus a 0.5\% margin, and the reserve
$\mathrm{SOC}_{\mathrm{res}}$.

\paragraph{Lagrangian optimization}
The planner selects $\lambda^\ast$ from a grid
$\Lambda = \{0, 0.05, \dots, 30, 10^3\}$ (15 values). For each $\lambda$ it predicts the
remaining trajectory at the forecast complexity level $\hat c_{\mathrm{fc}}$ (the mean of $\hat c$ over the mission so far):
the speed $v^\ast(\lambda)$ and energy per meter $e(\lambda)$ from \eqref{eq:schedule} with the
corrected power model $\hat P$, and the cell energy needed for the remaining distance $d$,
$\hat E_{\mathrm{cell}}(\lambda) = \hat\eta\,(1+\epsilon)\, e(\lambda)\, d$, where
$\epsilon = 5\%$ is a safety margin. The planner minimizes the predicted remaining cost
\begin{equation}
\begin{aligned}
  \lambda^\ast = \arg\min_{\lambda \in \Lambda}\; &
  w_T \frac{d}{v^\ast(\lambda)\, T_{\mathrm{ref}}}
  + w_E \frac{\hat\eta\, e(\lambda)\, d}{L}
  + w_R \min\!\big(1, \hat\phi(\lambda)\big)
  + w_D \frac{\widehat{\Delta\mathrm{SOH}}(\lambda)}{1 - \mathrm{SOH}_{\mathrm{EOL}}} \\
  \text{s.t.} \quad &
  \hat E_{\mathrm{cell}}(\lambda) \leq E_{\mathrm{av}}(\widehat{\mathrm{SOC}}, \widehat{\mathrm{SOH}}; \mathrm{SOC}_{\mathrm{f}}),
\end{aligned}
  \label{eq:lambda}
\end{equation}
where $\hat\phi(\lambda)$ is the predicted shortfall below the reserve, i.e.,
$\big[\hat E_{\mathrm{cell}}(\lambda) - E_{\mathrm{av}}(\cdot; \mathrm{SOC}_{\mathrm{res}})\big]^+$
converted to \ac{SOC} and normalized by $\mathrm{SOC}_{\mathrm{res}}$, and
$\widehat{\Delta\mathrm{SOH}}(\lambda)$ is the capacity loss \eqref{eq:dsoh} predicted for the
trajectory (see below).
The constraint discards prices whose trajectory is predicted to deplete the battery,
while the reserve term lets the planner use the reserve only where the time saved outweighs its penalty.
Because $E_{\mathrm{av}}$ in \eqref{eq:eav} depends on the estimated \ac{SOH}, the same
\ac{SOC} yields a higher energy price for an aged battery, which lowers the speed
and, through \eqref{eq:schedule}, the \ac{CPU} frequency.
If no $\lambda \in \Lambda$ is feasible, the requirement cannot be met with the remaining
charge, and the controller switches to a budget mode that keeps the energy within
$E_{\mathrm{av}}(\cdot; \mathrm{SOC}_{\mathrm{f}})$ and maximizes the achievable throughput
ratio $\min(1, q / q_{\mathrm{ref}})$, i.e., completion takes priority over \ac{QoS}.

\paragraph{Degradation prediction (DARO)}
DARO predicts the mission damage \eqref{eq:fd_mission} from the estimated battery state. At time $t$,
let $\hat E_{\mathrm{cell}}$ be the predicted cell energy for the remaining distance and
$\hat t_{\mathrm{r}} = d / v^\ast(\lambda)$ the remaining time. The predicted end \ac{SOC} is
\begin{equation}
  \widehat{\mathrm{SOC}}_{\mathrm{end}} = \widehat{\mathrm{SOC}}
    - \frac{\hat E_{\mathrm{cell}}}{Q_n\,\widehat{\mathrm{SOH}}\; V_{\mathrm{oc}}(\tilde\sigma, \widehat{\mathrm{SOH}})},
  \quad
  \tilde\sigma = \widehat{\mathrm{SOC}}
    - \frac{\hat E_{\mathrm{cell}}}{2\,Q_n\,\widehat{\mathrm{SOH}}\; V_{\mathrm{oc}}(\widehat{\mathrm{SOC}}, \widehat{\mathrm{SOH}})},
  \label{eq:soc_end_pred}
\end{equation}
i.e., the \ac{OCV} is evaluated at the midpoint $\tilde\sigma$ of the remaining discharge. With
$\hat\delta = \mathrm{SOC}_0 - \widehat{\mathrm{SOC}}_{\mathrm{end}}$,
$\hat{\bar\sigma} = \tfrac{1}{2}(\mathrm{SOC}_0 + \widehat{\mathrm{SOC}}_{\mathrm{end}})$, the mean measured
temperature $\bar T_{\mathrm{p}}$ over the elapsed time $t$, the current temperature $T_t$, and
$\hat{\bar T} = (t\,\bar T_{\mathrm{p}} + \hat t_{\mathrm{r}}\, T_t)/(t + \hat t_{\mathrm{r}})$, the
predicted damage is
\begin{equation}
\begin{aligned}
  \Delta\hat f_d = {} & \tfrac{1}{2} k_{\mathrm{cyc}}\, \widetilde{S}_{\hat\delta} S_{\hat{\bar\sigma}}
      \big(S_{\hat{\bar T}} + S_{T_{\mathrm{amb}}}\big) \\
  & + \Delta f_{\mathrm{cal}}\big(\mathrm{SOC}_0, \widehat{\mathrm{SOC}}, t, \bar T_{\mathrm{p}}\big)
    + \Delta f_{\mathrm{cal}}\big(\widehat{\mathrm{SOC}}, \widehat{\mathrm{SOC}}_{\mathrm{end}}, \hat t_{\mathrm{r}}, T_t\big) \\
  & + \Delta f_{\mathrm{cal}}\big(\widehat{\mathrm{SOC}}_{\mathrm{end}}, \mathrm{SOC}_0, \hat\delta / C_{\mathrm{ch}}, T_{\mathrm{amb}}\big),
\end{aligned}
  \label{eq:fd_pred}
\end{equation}
which replaces the measured trajectory in \eqref{eq:fd_mission} by its elapsed and predicted remaining
parts. The degradation cost used in \eqref{eq:lambda} is
\begin{equation}
  D(\hat E_{\mathrm{cell}}, \hat t_{\mathrm{r}})
  = w_D\, \frac{\widehat{\mathrm{SOH}}\,\big(1 - e^{-\Delta\hat f_d}\big)}{1 - \mathrm{SOH}_{\mathrm{EOL}}}.
  \label{eq:deg_cost}
\end{equation}
Evaluating $D$ for every speed and \ac{DVFS} level at every step is unnecessary, because one
step changes the remaining energy and time only marginally. The planner therefore linearizes $D$
at the trajectory of the selected price $\lambda^\ast$ by forward differences,
\begin{equation}
\begin{aligned}
  \mu_E &= \Big[\tfrac{D(\hat E_{\mathrm{cell}} + \Delta E,\, \hat t_{\mathrm{r}}) - D(\hat E_{\mathrm{cell}},\, \hat t_{\mathrm{r}})}{\Delta E}\Big]^+, \\
  \mu_t &= \Big[\tfrac{D(\hat E_{\mathrm{cell}},\, \hat t_{\mathrm{r}} + \Delta t_{\mathrm{r}}) - D(\hat E_{\mathrm{cell}},\, \hat t_{\mathrm{r}})}{\Delta t_{\mathrm{r}}}\Big]^+,
\end{aligned}
  \label{eq:deg_prices}
\end{equation}
with $\Delta E = \max(0.01\,\hat E_{\mathrm{cell}}, 1~\mathrm{J})$ and
$\Delta t_{\mathrm{r}} = \max(0.01\,\hat t_{\mathrm{r}}, 1~\mathrm{s})$. $\mu_E$ is the
marginal cycle aging per joule of cell energy (through the depth of discharge) and $\mu_t$ the
marginal calendar aging per second; both enter the cost per meter \eqref{eq:lagrangian}, so the
decision-making accounts for the degradation without evaluating the degradation model at every step.

\paragraph{Algorithm and complexity}
\begin{algorithm}[t]
\caption{DARO (Lag-RLS for $w_D = 0$)}
\label{alg:lag_rls}
\begin{algorithmic}[1]
\State Initialize $\boldsymbol\theta \gets \mathbf 1$, $\hat\eta \gets 1$, $\widehat{\mathrm{SOH}} \gets 1$, $\lambda^\ast \gets 0$, $\mu_E, \mu_t \gets 0$
\For{each control step $t$}
  \State Measure $v_t, f_t, q_t, V_t, I_t, \mathrm{SOC}^{\mathrm{cc}}_t, d$
  \State \emph{Estimator:} update $\widehat{\mathrm{SOH}}$, $\widehat{\mathrm{SOC}}$, $\hat\eta$, $\boldsymbol\theta$ (\ac{RLS}) and $\hat c$
  \If{$t \bmod T_p = 0$} \Comment{Lagrangian planner}
    \State Compute $E_{\mathrm{av}}$ at $\mathrm{SOC}_{\mathrm{f}}$ and $\mathrm{SOC}_{\mathrm{res}}$ by \eqref{eq:eav}
    \State Select $\lambda^\ast$ by \eqref{eq:lambda}; update $\mu_E, \mu_t$
    \State \textbf{if} no $\lambda$ is feasible \textbf{then} switch to budget mode
  \EndIf
  \State \emph{Decision-making:} $(v^\ast, f^\ast)$ by \eqref{eq:schedule}
  \State Apply speed probing during the first 90~s; send $(v^\ast, f^\ast)$
\EndFor
\end{algorithmic}
\end{algorithm}
Algorithm~\ref{alg:lag_rls} summarizes the controller. The decision-making evaluates
\eqref{eq:lagrangian} on the $|\mathcal{V}||\mathcal{F}|$ speed--frequency pairs; the expected
penalty \eqref{eq:qos_cost} depends only on $\hat c$ and is recomputed only when $\hat c$ changes.
The planner repeats the search for $|\Lambda|$ prices every $T_p$ steps, and the \ac{SOH} estimator
evaluates $|\mathcal{H}|$ voltage predictions per step. The cost per step is therefore
$\mathcal{O}(|\mathcal{V}||\mathcal{F}| + |\mathcal{H}| + |\Lambda||\mathcal{V}||\mathcal{F}|/T_p)$, and the mean
decision time in our Python implementation is about 19~ms, well below the 1~s control period.
